O gênero de jogos \textit{roguelike}, conhecido pela alta rejogabilidade
proporcionada pela \gls{PCG}, frequentemente apresenta desafios no que tange ao
controle fino e à criação de conteúdo com intenções específicas por parte dos
desenvolvedores. A manipulação de parâmetros complexos e pouco flexíveis em
algoritmos de \gls{PCG} tradicionais torna o processo de design pouco intuitivo.
Este trabalho explora uma abordagem para superar essa limitação, propondo um
sistema que utiliza \gls{LLM} para gerar mapas a partir de descrições textuais
em linguagem natural.  Reconhecendo as limitações inerentes das LLMs em
raciocínio espacial e geométrico, a pesquisa desenvolveu uma arquitetura híbrida
que desacopla responsabilidades: algoritmos determinísticos tradicionais
gerenciam a topologia estrutural e a conectividade da masmorra, enquanto a
\gls{LLM} atua como um "agente criativo"\ responsável pela geração semântica de
narrativas, inimigos e itens.  Para materializar esses elementos, foi
implementado um \textit{pipeline} que integra \textit{Structured Outputs} e a
técnica de \gls{RAG} sobre um banco vetorial curado de \textit{tiles},
garantindo a consistência visual e mitigando alucinações. Um jogo funcional foi
desenvolvido para interpretar os ativos gerados e renderizar níveis jogáveis. A
validação, realizada através de métricas de Reconstrução Semântica e Coerência
sobre três modelos distintos, demonstrou que modelos menores executados
localmente apresentam desempenho comparável a grandes modelos. Os resultados
confirmam a viabilidade técnica de associar linguagem natural a ativos de jogo
concretos, estabelecendo uma base sólida para sistemas de criação
\textit{"text-to-map"} que preservam a intenção do usuário.

\palavraschave{geração procedural de conteúdo; modelos de linguagem de larga
escala; jogos roguelike; geração de mapas.}